\documentclass{article}
\usepackage[utf8]{inputenc}
\usepackage[T1]{fontenc}
\usepackage[spanish]{babel}
\usepackage{amsmath, amssymb, amsthm}
\usepackage{geometry}
\geometry{margin=3cm}

\begin{document}

\begin{center}
  \Large\bfseries Demostración combinatoria del Teorema de Lucas
\end{center}

\bigskip

\noindent
\textbf{Enunciado.} Sea $p$ un número primo y sean $a,b\in\mathbb N$ con
\[
  a = a_0 + a_1p + a_2p^2 + \dots + a_kp^k,
  \qquad
  b = b_0 + b_1p + b_2p^2 + \dots + b_kp^k,
\]
donde cada dígito $a_i,b_i$ pertenece a $\{0,1,\dots,p-1\}$.  
El \emph{Teorema de Lucas} afirma
\[
  \binom{a}{b}
  \equiv
  \binom{a_0}{b_0}\binom{a_1}{b_1}\binom{a_2}{b_2}\cdots\binom{a_k}{b_k}
  \pmod{p}.
\]

A continuación se presenta una demostración puramente combinatoria,
sin recurrir a acciones de grupo ni a simetrías cíclicas, basada sólo en
el principio de adición, el principio de multiplicación y particiones
del conjunto que se está contando.

\section*{Interpretación combinatoria de {$\binom{a}{b}$}{(a b)}}

El coeficiente binomial $\binom{a}{b}$ cuenta los subconjuntos de tamaño
$b$ de un conjunto de $a$ elementos, que denotaremos
$S=\{1,2,\dots,a\}$.

\section*{2.\;Partición de {$S$}{S} en bloques según potencias de
{$p$}{p}}

Del desarrollo en base $p$ de $a$ se deduce la siguiente partición:

\begin{itemize}
  \item $a_1$ \emph{bloques completos} de tamaño $p$;
  \item un \emph{bloque incompleto} de tamaño $a_0$ (si $a_0=0$
        no existe este último bloque).
\end{itemize}

En símbolos:
\[
  a = a_1p + a_0, \qquad 0\le a_0 < p .
\]

\section*{3.\;Distribución de la elección de $b$ elementos}

Escribamos igualmente
$b = b_0 + b_1p$ con $0\le b_0<p$.
Sea
\[
  (x_1,x_2,\dots,x_{a_1},x_{\text{last}})
\quad
\bigl(0\le x_i\le p,\;0\le x_{\text{last}}\le a_0\bigr)
\]
una $($y sólo una$)$ forma de repartir los $b$ elementos entre los
bloques, sujeta a
\[
  x_1+\dots +x_{a_1}+x_{\text{last}} = b .
\]

\paragraph{Principio de multiplicación.}
Para una tupla fija se tiene
\[
  N(x_1,\dots,x_{a_1},x_{\text{last}})
  = \binom{p}{x_1}\binom{p}{x_2}\cdots\binom{p}{x_{a_1}}
    \binom{a_0}{x_{\text{last}}}.
\]

\paragraph{Principio de adición.}
Sumando sobre todas las tuplas viables,
\[
  \binom{a}{b}
  =\!\!\!
    \sum_{\substack{x_1+\dots+x_{a_1}+x_{\text{last}}=b\\
                    0\le x_i\le p,\,0\le x_{\text{last}}\le a_0}}
      \binom{p}{x_1}\binom{p}{x_2}\cdots\binom{p}{x_{a_1}}
      \binom{a_0}{x_{\text{last}}}.
\]

\section*{4.\;Cálculo módulo {$p$}{p}}

Para $1\le x\le p-1$ se cumple
$\displaystyle\binom{p}{x}\equiv 0\pmod{p}$,
mientras que
$\binom{p}{0}=\binom{p}{p}=1$.

Por tanto, \emph{sólo} aportan al valor de la suma módulo $p$
aquellas tuplas en las que cada $x_i$ vale $0$ o $p$.
Sea $t$ el número de índices con $x_i=p$; entonces
\[
  tp + x_{\text{last}} = b \quad\Longrightarrow\quad
  x_{\text{last}} = b-tp.
\]

La desigualdad $0\le x_{\text{last}}\le a_0$
junto con $0\le b_0 <p$ fuerza necesariamente
$t=b_1$ y $x_{\text{last}}=b_0$.
Además $b_0\le a_0$ es condición necesaria y suficiente
para que existan selecciones (si falla, ambos lados de la congruencia
son $0$ y Lucas se verifica trivialmente).

\medskip
\noindent
\textbf{Cuenta de las configuraciones supervivientes.}
\begin{itemize}
  \item Elegir los $b_1$ bloques completos que se toman enteros:
        $\binom{a_1}{b_1}$ maneras.
  \item En cada bloque seleccionado la elección es única
        (se toman sus $p$ elementos).
  \item En el bloque incompleto se eligen $b_0$ de sus $a_0$ elementos:
        $\binom{a_0}{b_0}$ maneras.
\end{itemize}
De donde
\[
  \binom{a}{b}\equiv \binom{a_1}{b_1}\binom{a_0}{b_0}\pmod{p}.
\]

\section*{5.\;Iteración sobre los dígitos superiores}

El argumento se aplica recursivamente al binomio
$\binom{a_1}{b_1}$: se agrupan ahora los $a_1$ bloques completos
en super‑bloques de tamaño $p^2$, regidos por los dígitos $a_2,b_2$,
y se repite el razonamiento.
Encadenando todos los pasos (hasta alcanzar $a_k,b_k$) se obtiene
\[
  \boxed{\;
    \displaystyle
    \binom{a}{b}
    \equiv
    \prod_{i=0}^k \binom{a_i}{b_i}
    \pmod{p}}\;,
\]
que es exactamente el Teorema de Lucas.

\section*{6.\;Ejemplo ilustrativo}

Para $p=5$, $a=12$, $b=7$ (en base $5$, $a=(2\,2)_5$, $b=(1\,2)_5$)
Lucas predice
\[
  \binom{12}{7}\equiv
  \binom{a_0}{b_0}\binom{a_1}{b_1}
  =\binom{2}{2}\binom{2}{1}=2\pmod{5}.
\]
Un desglose de las particiones de
$\{1,\dots,12\}$
en dos bloques completos de $5$ y uno incompleto de $2$
muestra que sólo dos subconjuntos
(\,$\{1,\dots,5,11,12\}$ y $\{6,\dots,10,11,12\}$\,)
producen contribución no nula módulo $5$, confirmando
el resultado $\binom{12}{7}\equiv 2\pmod5$.

\bigskip
\hrule
\bigskip

\noindent
\textbf{Conclusión.}  
La clave fue aislar, mediante particiones
controladas por las potencias de $p$, las únicas selecciones
que sobreviven módulo $p$:
aquellas donde cada bloque completo se toma entero o se deja intacto.
Eso traduce la elección global en elecciones dígito a dígito, dando lugar
al producto $\prod_i \binom{a_i}{b_i}$ que establece el teorema.

% --- Fragmento a insertar tras la última subsección ---
\subsection*{Conteo de elecciones distribuidas en bloques (Principio de adición y multiplicación)}

Consideremos cuántas maneras hay de elegir un subconjunto de $b$ elementos de $S$ \emph{distribuyendo} esos $b$ elementos entre los bloques definidos.  
Escribamos $b$ en base $p$ como $b = b_1 p + b_0$ con $0 \le b_0 < p$.  
Al seleccionar $b$ elementos de $S$, fijemos que elegimos

\begin{itemize}
  \item $x_i$ elementos del bloque $i$ (para cada bloque completo $i = 1,2,\dots,a_1$), donde cada bloque completo tiene $p$ elementos y, por tanto, $0 \le x_i \le p$;
  \item $x_{\text{last}}$ elementos del bloque incompleto, con $0 \le x_{\text{last}} \le a_0$.
\end{itemize}

Estas cantidades deben satisfacer la restricción
\[
  x_1 + x_2 + \cdots + x_{a_1} + x_{\text{last}} \;=\; b,
\]
pues en total se eligen $b$ objetos.

\paragraph{Principio de adición (partición de casos).}
Cada elección de $b$ elementos se clasifica según la tupla
$(x_1,\dots,x_{a_1},x_{\text{last}})$ que produce.  
Diferentes tuplas definen subconjuntos disjuntos de selecciones
y la unión de todas ellas cubre la totalidad de los $b$‑subconjuntos.
Así, por el principio de adición, el número total de formas (esto es, $\binom{a}{b}$) es la \emph{suma} de las cantidades correspondientes a cada tupla válida.

\paragraph{Principio de multiplicación (conteo dentro de cada categoría).}
Fijemos ahora una tupla válida $(x_1,\dots,x_{a_1},x_{\text{last}})$.
Las elecciones en bloques distintos son independientes, de modo que:

\begin{itemize}
  \item del bloque $1$ podemos elegir $x_1$ elementos de $\binom{p}{x_1}$ formas;
  \item del bloque $2$, $\binom{p}{x_2}$ formas, y así sucesivamente;
  \item del bloque $a_1$, $\binom{p}{x_{a_1}}$ formas;
  \item del bloque incompleto, $\binom{a_0}{x_{\text{last}}}$ formas.
\end{itemize}

Por el principio de multiplicación, el número total de subconjuntos
asociados a esa tupla es
\[
  \binom{p}{x_1}\,\binom{p}{x_2}\,\cdots\,
  \binom{p}{x_{a_1}}\,
  \binom{a_0}{x_{\text{last}}}.
\]

\noindent Sumando sobre \emph{todas} las tuplas que satisfacen la restricción, obtenemos
\[
  \binom{a}{b}
  \;=\;
  \sum_{\substack{
      x_1+\dots+x_{a_1}+x_{\text{last}}\,=\,b\\
      0 \le x_i \le p,\;
      0 \le x_{\text{last}} \le a_0
  }}
  \binom{p}{x_1}\,
  \binom{p}{x_2}\,\cdots\,
  \binom{p}{x_{a_1}}\,
  \binom{a_0}{x_{\text{last}}}.
\]

Esta identidad cuenta exhaustivamente todas las maneras de distribuir
$b$ elementos entre los bloques, \emph{clasificándolas} por cuántos
se toman de cada bloque y \emph{multiplicando} las elecciones
independientes dentro de cada bloque.


% --- Fragmento a insertar tras la sección anterior ---
\subsection*{Análisis módulo $p$: identificando contribuciones no nulas}

Examinemos ahora la expresión obtenida \emph{módulo $p$}, aprovechando que $p$ es primo.  
El hecho clave es:

\begin{quote}
Para un bloque completo de $p$ elementos se verifica
\[
  \binom{p}{x}\equiv 0 \pmod{p}
  \quad\text{si}\quad 1\le x\le p-1,
  \qquad
  \binom{p}{0}=\binom{p}{p}=1.
\]
\end{quote}

El motivo es combinatorio: cualquier subconjunto no trivial de un conjunto de
$p$ elementos se puede rotar cíclicamente $p$ veces, y las $p$ rotaciones son
distintas cuando $p$ es primo; así, el total de tales subconjuntos es múltiplo
de $p$.

Considérese un término de la suma
\[
  T_{(x_1,\dots,x_{a_1},x_{\text{last}})}
  \;=\;
  \binom{p}{x_1}\cdots\binom{p}{x_{a_1}}\,
  \binom{a_0}{x_{\text{last}}}.
\]
El producto $T_{(\cdot)}$ \emph{no} será divisible por $p$ únicamente cuando
cada factor $\binom{p}{x_i}$ carezca de $p$ como divisor, es decir,
\[
  x_i\in\{0,p\}\quad\text{para todo bloque completo }i.
\]
El factor $\binom{a_0}{x_{\text{last}}}$ nunca introduce un divisor $p$
porque $a_0<p$.

\medskip
\noindent
Sean, entonces, $x_i\in\{0,p\}$ y escribamos
\[
  t=\#\bigl\{\,i\;|\;x_i=p\bigr\},
  \qquad
  x_{\text{last}}=b-tp.
\]
La ecuación de conteo total
\[
  x_1+\dots+x_{a_1}+x_{\text{last}}=b
\]
se transforma en
\[
  t\,p + x_{\text{last}} = b,
  \qquad
  0\le x_{\text{last}}\le a_0.
\]

\paragraph{Determinación de $t$.}
Escribiendo $b=b_1p+b_0$ con $0\le b_0<p$:

\begin{itemize}
  \item Si $t<b_1$ entonces
        $x_{\text{last}} = b - tp > b - b_1p = b_0$,
        excediendo el límite permitido $a_0\ge b_0$.
  \item Si $t>b_1$ entonces $tp\ge(b_1+1)p>b$, imposible.
\end{itemize}

\noindent
La única posibilidad es
\[
  t=b_1,\qquad x_{\text{last}}=b_0.
\]
Es decir, \emph{exactamente} $b_1$ bloques completos se toman
\emph{enteros} y del bloque incompleto se eligen $b_0$ elementos.
La condición $b_0\le a_0$ (y, más generalmente, $b_i\le a_i$ para
todos los dígitos) garantiza la existencia de la selección.

Los términos sobrevivientes módulo $p$ corresponden, pues, a elegir un
subconjunto que:
\begin{enumerate}
  \item incluye \textbf{todos} los elementos de exactamente $b_1$ bloques
        completos,
  \item excluye los elementos de los $a_1-b_1$ bloques completos restantes,
  \item contiene $b_0$ elementos del bloque incompleto.
\end{enumerate}

\medskip
\noindent
El número de tales selecciones es
\[
  \binom{a_1}{b_1}\,\binom{a_0}{b_0}
  \;\equiv\;
  \binom{a_1}{b_1}\,\binom{a_0}{b_0}
  \pmod{p},
\]
donde $\binom{a_1}{b_1}$ cuenta la elección de los $b_1$ bloques completos
que se toman enteros.  
Puesto que $a_1$ y $b_1$ son, respectivamente, los dígitos de $a$ y $b$
en la posición $p^1$, obtenemos precisamente el factor
$\binom{a_1}{b_1}\binom{a_0}{b_0}$ previsto por el enunciado de Lucas.

Repitiendo el mismo argumento en los niveles $p^2,p^3,\dots,p^k$ se
obtiene inductivamente
\[
  \binom{a}{b}
  \;\equiv\;
  \prod_{i=0}^{k}
  \binom{a_i}{b_i}
  \pmod{p},
\]
completando la demostración.


\end{document}
